\begin{abstract}
Multi-agent collaborative perception could significantly upgrade the perception performance by enabling agents to share complementary information with each other through communication. It inevitably results in a fundamental trade-off between perception performance and communication bandwidth. To tackle this bottleneck issue, we propose a spatial confidence map, which reflects the spatial heterogeneity of perceptual information. It empowers agents to only share spatially sparse, yet perceptually critical information, contributing to where to communicate. Based on this novel spatial confidence map, we propose \texttt{Where2comm}, a communication-efficient collaborative perception framework. \texttt{Where2comm} has two distinct advantages: i) it considers pragmatic compression and uses less communication to achieve higher perception performance by focusing on perceptually critical areas; and ii) it can handle varying communication bandwidth by dynamically adjusting spatial areas involved in communication. To evaluate~\texttt{Where2comm}, we consider 3D object detection in both real-world and simulation scenarios with two modalities (camera/LiDAR) and two agent types (cars/drones) on four datasets: OPV2V, V2X-Sim, DAIR-V2X, and our original CoPerception-UAVs.~\texttt{Where2comm} consistently outperforms previous methods; for example, it achieves more than $100,000 \times$ lower communication volume and still outperforms DiscoNet and V2X-ViT on OPV2V. Our code is available at~\url{https://github.com/MediaBrain-SJTU/where2comm}.
\vspace{-2mm}
% by strategically allocating communication budges on spatial areas based on their perceptual criticalness
% With agents sharing complementary information with each other by communication, multi-agent collaborative perception systems have demonstrated breakthrough performances in perception tasks. Yet, real-world communication among agents can be expensive if the bandwidth is limited. For communication efficiency, we introduce a spatial confidence map that reflects the perceptually critical levels of each agent' perceived information on the spatial dimension.
% % we identify a spatial-confidence-aware strategy to let agents selectively share and use the more critical information in a communication system. Each agent generates a spatial confidence map that reflects the perceptually critical levels of its perceived information on the spatial dimension. 
% Based on the confidence map, the agent identifies what area of information (where) is more critical  to be shared than others, which maximizes the use of bandwidth. Based on this strategy, we propose \texttt{Where2comm}, a novel multi-round multi-agent collaborative perception framework. \texttt{Where2comm} leverages the spatial confidence map to i) avoid bandwidth waste by letting agents only pack the critical information in the messages; ii) promote targeted communication by enabling each agent to send spatial-specific information requests; iii) build a more interpretable communication graph; iv) support confidence-aware message fusion. To comprehensively evaluate \texttt{Where2comm}, we conduct experiments on three datasets: OPV2V~\cite{OPV2V}, V2X-Sim~\cite{V2XSim}, and our original dataset CoPerception-UAVs, covering two types of agents (cars and drones) and sensors (LiDAR and cameras). We achieve more than $5,000$ times less communication volume, and still outperforms the state-of-the-art methods for both AP@0.50 and 0.70 on OPV2V and CoPerception-UAVs.

%Multi-agent collaborative perception systems have demonstrated breakthrough performances in many perception tasks by enabling agents to share complementary information with each other through communication. However, in real world applications, the communication among agents can be expensive when the bandwidth is limited. Considering the bandwidth and communication efficiency, each agent in the system should selectively pick the information that is worth sharing to join the communication. Following this thread of thought, we come up with a spatial-confidence-aware strategy. It enables a spatial confidence map to reflect the perceptually critical levels for each agent's perceived information on the spatial dimension. Each agent's confidence map guides it to clearly identify what area of information is more critical than others and to only share the critical information to make every bit of the bandwidth count. Based on this strategy, we propose \texttt{Where2comm}, a novel multi-round multi-agent collaborative perception framework. \texttt{Where2comm} leverages the spatial confidence map to i) avoid bandwidth waste by letting agents only pack the critical information in the messages; ii) promote targeted communication by enabling each agent to send spatial-specific information requests; iii) build a more interpretable communication graph; iv) support confidence-aware message fusion. To comprehensively evaluate \texttt{Where2comm}, we conduct experiments on three datasets: OPV2V~\cite{OPV2V}, V2X-Sim~\cite{V2XSim}, and our original dataset CoPerception-UAVs, covering two types of agents (cars and drones) and two types of sensors (LiDAR and cameras). We achieve more than $5,000$ times less communication volume and still outperforms the state-of-the-art methods for both AP@0.50 and 0.70 on OPV2V and CoPerception-UAVs.

% Collaborative perception employs cross-agent communication to provide coverage and robustness benefits for single-agent perception; however, perception data consumes huge communication bandwidth, making it difficult to be practically deployed in real-world applications. 
% To promote better performance-bandwidth trade-off, we propose a novel spatial-confidence-aware communication strategy, which enables each agent to flexibly decide where to communicate by identifying the perceptual critical level at each location. Each agent could provide more critical information for other agents, and requesting more targeted information from other agents in multi-round communication. 
% Following this strategy, we propose \texttt{Where2comm}, a novel multi-round and multi-agent collaborative perception framework. \texttt{Where2comm} leverages the spatial confidence map to decide where to communicate via novel message packing, and who to communicate via novel communication graph construction; and whose information to aggregate via effective prior-based message fusion.
% To comprehensively evaluate \texttt{Where2comm}, we conduct experiments on three datasets: OPV2V~\cite{OPV2V}, V2X-Sim~\cite{V2XSim}, and our collected dataset CoPerception-UAVs, covering two types of agents (cars and drones) and two types of sensors (LiDAR and cameras). We achieve more than $5,000$ times less communication volume and still outperforms the state-of-the-art methods for both AP@0.50 and 0.70 on OPV2V and CoPerception-UAVs.


% Multi-agent collaborative perception can essentially solve the inevitable physics limitation of single-agent perception such as occlusion and long-range issues. However, existing approaches place far more than practical demands on communication systems. To achieve a feasible bandwidth-performance trade-off, we propose a novel spatial-aware collaboration mechanism, Where2comm, where agents generate confidence and request maps based on their own views to decide where to communicate and how to pack messages. Based on the guidance of the proposed confidence and request map, we design a multi-round and dynamic bandwidth adaptation communication strategy, to leverage a targeted, precise and efficient collaboration. In order to efficiently integrate limited information, we propose a confidence-aware multi-head attention for information fusion purposes to reach better perception performance. We conduct extensive experiments on multiple datasets on different views and modalities to prove that Where2comm achieves states of the art performance and bandwidth-efficiency.

% We propose a novel fine-grained spatial-aware communication strategy, where each agent can decide where to communicate and pack messages only related to the most perceptually critical spatial areas. This strategy not only enables more precise support for other agents, but also more targeted request from other agents in multi-round communication. We propose~\texttt{Where2comm}, a novel collaborative perception framework based on the spatial-aware communication strategy. With the guidance of the proposed spatial confidence map,~\texttt{Where2comm} leverages novel message packing and communication graph learning to  achieve lower communication bandwidth, and adopts confidence-aware multi-head attention to reach better perception performance. We conduct extensive experiments to validate~\texttt{Where2comm} achieves state-of-the-art performance-bandwidth trade-off on multiple challenging datasets across views and modalities. 
\end{abstract}